\subsection{Bloom-Filter Duplicate Detection}

We evaluated a Bloom filter~\cite{bloom1970} as a memory-efficient way to detect
duplicate events in the replay stream. The experiment is standalone: it uses the
event identifiers from the 18{,}000-event replay input but does not run through Kafka.

\subsubsection{Method}
The replay input contains 18{,}000 unique event identifiers, so duplicates were
injected synthetically. In each trial, 1{,}800 identifiers (10\%) were copied and
each copy was placed at a random later position, giving a stream of 19{,}800
events. Filters were sized for 18{,}000 distinct items at target false-positive
rates of 5\%, 1\% and 0.1\%, using $m=-n\ln p/(\ln 2)^2$ bits and
$k=(m/n)\ln 2$ hash functions. Hash positions were derived from SHA-256 using double hashing~\cite{kirsch2008less}, with a per-trial salt. Each configuration was run for five trials, which differ in which
events were duplicated, where the copies were placed, and the hash salt. An exact
in-memory set served as ground truth. A flagged event was a false positive if it
had not been seen before, and a missed duplicate if it had been seen but not flagged.

\subsubsection{Results}
Table~\ref{tab:bloom} summarises the results. All 1{,}800 injected duplicates were
detected in every trial, with no misses, as expected for a Bloom filter. The
measured false-positive rate was lower than the target in each case. The target
rate applies to a filter that already holds all items, whereas in the stream the
filter fills gradually, so most lookups occur when it is less full. The
corresponding full-filter rates were 5.03\%, 1.00\% and 0.10\%, which matches the
targets and confirms the filters were sized correctly.

\begin{table}[t]
\caption{Bloom-filter results (mean of five trials; 18{,}000 distinct events, 1{,}800 injected duplicates)}
\label{tab:bloom}
\centering
\begin{tabular}{lccc}
\hline
Target FP rate & 5\% & 1\% & 0.1\% \\
\hline
Hash functions & 4 & 7 & 10 \\
Filter size (KiB) & 13.7 & 21.1 & 31.6 \\
Duplicates detected & 1{,}800 & 1{,}800 & 1{,}800 \\
Duplicates missed & 0 & 0 & 0 \\
False positives & 222.2 & 29.0 & 2.2 \\
Measured FP rate & 1.23\% & 0.16\% & 0.012\% \\
\hline
\end{tabular}
\end{table}

For comparison, an exact Python set of the same identifiers used about 1.78~MiB,
so the filters were roughly 133, 87 and 58 times smaller. This figure is specific to
our Python implementation.

\subsubsection{Limitations}
The duplicates are synthetic because the replay input has no real duplicates, so
the results describe filter behaviour rather than duplicate rates in the source
data. The experiment was not integrated into the Kafka or Spark pipeline. A standard Bloom filter also fills up as events accumulate; Stable Bloom Filters~\cite{deng2006stable} address this for unbounded streams, and we did not evaluate them.
